Para pequeñas deformaciones:
\[
\mat{\varepsilon}
=\frac12\left[\nabla\vect u+(\nabla\vect u)^T\right].
\]

En una frontera donde imponemos desplazamientos:
\[
\vect u=\bar{\vect u}.
\]

\safesection{Comportamiento}
\[
\mat{\sigma}=\mathsf{C}:\mat{\varepsilon}.
\]

\safesection{La cadena fundamental}
\flowbox{$\vect u\ \xrightarrow{\text{cinemática}}\ \mat{\varepsilon}\ \xrightarrow{\text{material}}\ \mat{\sigma}\ \xrightarrow{\text{equilibrio}}\ \text{cargas}$}

Este esquema es fundamental porque reaparecerá más adelante en la asignatura.

\begin{tip}
Cuando no sepas ``qué ecuación falta'', pregúntate a qué familia pertenece:
\[
\text{equilibrio},\qquad\text{compatibilidad},\qquad\text{constitutiva}.
\]
Casi siempre localizarás así el problema.
\end{tip}

% ============================================================
